\section{Обоснование выбора технологий}
\label{sec:tech-choice}

Выбор технологического стека выполнялся с учётом подходов к проектированию информационных систем и итерационной разработке программного обеспечения~\cite{vendrov2006,gagarina2020,jacobson2002}. Для каждого слоя системы оценивались назначение компонента, характер нагрузки, требования к сопровождению и возможность независимого развёртывания.

\subsection{Язык программирования серверной части}

Для реализации основных серверных сервисов (аутентификация, маршруты, обработка
фотографий) выбран язык программирования Rust. Выбор связан с тем, что эти
компоненты обрабатывают сетевые запросы, выполняют сериализацию данных,
работают с изображениями и должны сохранять предсказуемое потребление ресурсов.
Основные аргументы:

\begin{itemize}
  \item компиляция в нативный код без сборщика мусора позволяет уменьшить
    непредсказуемые паузы выполнения, что важно для обработки изображений и
    длительных WebSocket-соединений~\cite{blandy2021};
  \item система владения и проверки заимствований на этапе компиляции исключает целые классы ошибок: обращение к освобождённой памяти, гонки данных, переполнение буфера~\cite{rustbook};
  \item экосистема Tokio предоставляет асинхронную среду выполнения для сетевых
    сервисов с большим числом одновременных операций ввода-вывода;
  \item фреймворк Axum обеспечивает эргономичный программный интерфейс для создания HTTP-сервисов, библиотека sqlx~--- компилируемые SQL-запросы с проверкой типов.
\end{itemize}

Для вспомогательных сервисов (кэширование тайлов, проксирование) выбран язык Go~--- прагматичный выбор для задач с интенсивным вводом-выводом с простым конкурентным программированием через горутины~\cite{donovan2015}.

\subsection{Клиентская часть}

Для разработки клиентской части выбран фреймворк React~19~\cite{banks2020} с языком TypeScript. Обоснование:
\begin{itemize}
  \item развитая экосистема компонентов для работы с картами (React Leaflet);
  \item строгая типизация TypeScript снижает количество ошибок на этапе разработки;
  \item поддержка режима прогрессивного веб-приложения для офлайн-работы;
  \item возможность обёртки в десктоп-приложение через Tauri.
\end{itemize}

\subsection{Система управления базами данных}

Для хранения данных выбрана СУБД PostgreSQL~16~\cite{silberschatz2019}. В
контексте данной системы важны следующие свойства:
\begin{itemize}
  \item поддержка типа JSONB для хранения гибких структур (точки маршрута с метаданными);
  \item расширение PostGIS для геопространственных запросов;
  \item транзакционность и развитые механизмы индексации для пользовательских и
    социальных данных;
  \item богатый набор индексов (B-tree, GIN для JSONB, GiST для геоданных).
\end{itemize}

\subsection{Архитектурный подход: микросервисы}

Для системы выбрана микросервисная архитектура вместо монолитной~\cite{newman2021,burns2018,tanenbaum2003}.
Такой выбор является инженерным компромиссом: он повышает сложность
развёртывания и сопровождения, но позволяет разделить компоненты с разным
характером нагрузки и разными требованиями к отказоустойчивости. Обоснование:
\begin{itemize}
  \item сервис обработки фотографий выполняет более длительные и ресурсоёмкие
    операции, чем сервис аутентификации или операции чтения маршрутов, поэтому
    его целесообразно масштабировать и перезапускать отдельно;
  \item сбой или задержка обработки фотографий не должны блокировать вход
    пользователя в систему и просмотр уже созданных маршрутов;
  \item сервисы тайлов и кеширования имеют вспомогательную сетевую нагрузку и
    могут развиваться независимо от доменной логики маршрутов;
  \item независимое развёртывание уменьшает риск обновления: изменение
    обработчика фотографий не требует пересборки сервиса аутентификации.
\end{itemize}

Такой подход согласуется с целью архитектурного проектирования, сформулированной
Р.~Мартином: <<уменьшить человеческие трудозатраты на создание и сопровождение
системы>>~\cite{martin2018}. Кроме того, выбранная архитектура учитывает
изменяемость требований: М.~Клеппман подчёркивает, что <<системные требования
навсегда останутся неизменными>> с вероятностью, стремящейся к нулю~\cite{kleppmann2017}.

\subsection{Брокер сообщений}

Для асинхронного взаимодействия микросервисов выбран NATS JetStream. Сравнение с альтернативами:
\begin{itemize}
  \item NATS против RabbitMQ. NATS проще в настройке и развёртывании, имеет минимальные требования к ресурсам (один исполняемый файл). RabbitMQ требует среду выполнения Erlang и более сложен в администрировании.
  \item NATS против Apache Kafka. Kafka ориентирован на обработку больших объёмов
    потоковых данных и избыточен для задач системы. NATS JetStream обеспечивает
    гарантию доставки <<не менее одного раза>> при меньшей эксплуатационной
    сложности для выбранного сценария.
\end{itemize}

\subsection{Сводная таблица сравнения технологий}

Результаты сравнительного анализа технологий по ключевым критериям представлены в таблицах~\ref{tab:tech-backend}, \ref{tab:tech-frontend} и~\ref{tab:tech-db}. Оценка ведётся по шкале 1--3 (1~--- неудовлетворительно, 2~--- приемлемо, 3~--- отлично). Удельный вес критерия отражает его значимость для данной системы.

\begin{spacing}{1.0}
\small
\begin{longtable}{|p{4.2cm}|c|c|c|c|c|}
\caption{Сравнение языков программирования серверной части}
\label{tab:tech-backend} \\
\hline
\textbf{Критерий} & \textbf{Вес} & \textbf{Rust} & \textbf{Go} & \textbf{Node.js} & \textbf{Java} \\
\hline
\endfirsthead
\multicolumn{6}{r}{\small Продолжение таблицы~\ref{tab:tech-backend}} \\
\hline
\endhead
\longtablecontinued{6}{tab:tech-backend}
\endfoot
\longtableended{6}{tab:tech-backend}
\endlastfoot
Производительность & 0{,}25 & 3 & 3 & 2 & 2 \\
\hline
Безопасность памяти & 0{,}20 & 3 & 2 & 1 & 2 \\
\hline
Асинхронность и конкурентность & 0{,}20 & 3 & 3 & 2 & 2 \\
\hline
Экосистема и библиотеки & 0{,}15 & 2 & 3 & 3 & 3 \\
\hline
Размер Docker-образа & 0{,}10 & 3 & 3 & 1 & 1 \\
\hline
Скорость разработки & 0{,}10 & 2 & 3 & 3 & 2 \\
\hline
\textbf{Взвешенная сумма} & --- & \textbf{2{,}80} & \textbf{2{,}80} & \textbf{1{,}90} & \textbf{2{,}10} \\
\hline
\end{longtable}
\end{spacing}

Rust и Go показали одинаковый взвешенный результат. Rust выбран для вычислительно-интенсивных сервисов (аутентификация, маршруты, обработка фотографий) из-за преимуществ в безопасности памяти и предсказуемой задержке. Go выбран для вспомогательных сервисов (кэш тайлов, прокси).

\begin{spacing}{1.0}
\small
\begin{longtable}{|p{4.2cm}|c|c|c|c|}
\caption{Сравнение фреймворков клиентской части}
\label{tab:tech-frontend} \\
\hline
\textbf{Критерий} & \textbf{Вес} & \textbf{React} & \textbf{Vue} & \textbf{Angular} \\
\hline
\endfirsthead
\multicolumn{5}{r}{\small Продолжение таблицы~\ref{tab:tech-frontend}} \\
\hline
\endhead
\longtablecontinued{5}{tab:tech-frontend}
\endfoot
\longtableended{5}{tab:tech-frontend}
\endlastfoot
Зрелость экосистемы & 0{,}25 & 3 & 2 & 3 \\
\hline
Поддержка TypeScript & 0{,}20 & 3 & 3 & 3 \\
\hline
Библиотеки для карт & 0{,}20 & 3 & 2 & 1 \\
\hline
Прогрессивное веб-приложение и десктоп-обёртка & 0{,}20 & 3 & 2 & 2 \\
\hline
Сообщество и документация & 0{,}15 & 3 & 3 & 2 \\
\hline
\textbf{Взвешенная сумма} & --- & \textbf{3{,}00} & \textbf{2{,}35} & \textbf{2{,}20} \\
\hline
\end{longtable}
\end{spacing}

React выбран как наиболее зрелый фреймворк с богатой экосистемой библиотек для интерактивных карт (React Leaflet) и поддержкой десктоп-обёртки Tauri.

\begin{spacing}{1.0}
\small
\begin{longtable}{|p{4.2cm}|c|c|c|c|}
\caption{Сравнение систем управления базами данных}
\label{tab:tech-db} \\
\hline
\textbf{Критерий} & \textbf{Вес} & \textbf{PostgreSQL} & \textbf{MongoDB} & \textbf{MySQL} \\
\hline
\endfirsthead
\multicolumn{5}{r}{\small Продолжение таблицы~\ref{tab:tech-db}} \\
\hline
\endhead
\longtablecontinued{5}{tab:tech-db}
\endfoot
\longtableended{5}{tab:tech-db}
\endlastfoot
JSONB для гибких структур & 0{,}25 & 3 & 3 & 1 \\
\hline
Геопространственные индексы & 0{,}25 & 3 & 2 & 1 \\
\hline
ACID-транзакции & 0{,}20 & 3 & 2 & 3 \\
\hline
Масштабируемость & 0{,}15 & 3 & 3 & 2 \\
\hline
Надёжность и зрелость & 0{,}15 & 3 & 2 & 3 \\
\hline
\textbf{Взвешенная сумма} & --- & \textbf{3{,}00} & \textbf{2{,}45} & \textbf{1{,}80} \\
\hline
\end{longtable}
\end{spacing}

PostgreSQL~\cite{silberschatz2019} выбран как единственная СУБД, обеспечивающая одновременно поддержку JSONB для хранения гибких структур точек маршрута и расширение PostGIS для геопространственных индексов и запросов.


\section{Обзор алгоритмов}
\label{sec:algorithms}

В рамках работы под алгоритмами, протоколами и форматами понимаются не сами внешние сервисы, а реализованные в системе процедуры обработки данных и обмена между компонентами. К таким элементам относятся: расчёт длины маршрута по координатам, оценка времени прохождения, классификация сложности маршрута, асинхронная обработка фотографий, фильтрация GPS-точек мобильного клиента, а также импорт и экспорт маршрутов в форматах GeoJSON, GPX и KML. При описании алгоритмической части учитывались общие подходы к анализу вычислительных процедур и обработке структурированных данных~\cite{skiena2011}.

\subsection{Формула Гаверсинуса}

Для расчёта расстояния между двумя точками маршрута по их географическим координатам используется формула Гаверсинуса (Haversine formula). Эта формула вычисляет расстояние по дуге большого круга между двумя точками на сфере по их широте и долготе.

Для двух точек с координатами $(\varphi_1, \lambda_1)$ и $(\varphi_2, \lambda_2)$, где $\varphi$~--- широта, $\lambda$~--- долгота (в радианах):

\begin{equation}
  a = \sin^2\!\left(\frac{\Delta\varphi}{2}\right) + \cos\varphi_1 \cdot \cos\varphi_2 \cdot \sin^2\!\left(\frac{\Delta\lambda}{2}\right)
\end{equation}

\begin{equation}
  c = 2 \cdot \mathrm{atan2}\!\left(\sqrt{a},\;\sqrt{1 - a}\right)
\end{equation}

\begin{equation}
  d = R \cdot c
\end{equation}

где $\Delta\varphi = \varphi_2 - \varphi_1$, $\Delta\lambda = \lambda_2 - \lambda_1$, $R = 6{,}371$~км~--- средний радиус Земли.

Общее расстояние маршрута вычисляется как сумма расстояний между последовательными точками:
\begin{equation}
  D = \sum_{i=1}^{n-1} d(\text{точка}_i, \text{точка}_{i+1})
\end{equation}

Формула Гаверсинуса обеспечивает точность порядка 0,5\% для расстояний менее 10\,000~км, что достаточно для туристических маршрутов.

\subsection{Формула Нейсмита}

Для оценки времени прохождения пешего маршрута используется формула Нейсмита (Naismith's rule), разработанная шотландским альпинистом Уильямом Нейсмитом в 1892~году. Правило устанавливает, что пешеход проходит 5~км в час по ровной местности, и дополнительно требуется 1~час на каждые 600~м набора высоты:

\begin{equation}
  T = \frac{D}{V} + \frac{H^{+}}{600}
\end{equation}

где $T$~--- время в часах, $D$~--- горизонтальное расстояние в километрах, $V = 5$~км/ч~--- скорость на ровной местности, $H^{+}$~--- суммарный набор высоты в метрах.

Набор высоты вычисляется по данным высот точек маршрута, полученным через Open-Meteo Elevation API:
\begin{equation}
  H^{+} = \sum_{i=1}^{n-1} \max(0,\; h_{i+1} - h_i)
\end{equation}

где $h_i$~--- высота $i$-й точки маршрута над уровнем моря.

\subsection{Алгоритм классификации сложности маршрута}

Для автоматической классификации сложности маршрута разработан алгоритм на основе комбинации метрик расстояния и набора высоты. Алгоритм присваивает маршруту одну из трёх категорий сложности: <<Лёгкий>>, <<Средний>>, <<Сложный>>.

Входные параметры: общее расстояние $D$ (км), суммарный набор высоты $H^{+}$ (м). Каждый параметр оценивается по балльной шкале:

\begin{itemize}
  \item расстояние: $D < 10$~км~--- 1 балл, $10 \leq D < 20$~км~--- 2 балла, $D \geq 20$~км~--- 3 балла;
  \item набор высоты: $H^{+} < 500$~м~--- 1 балл, $500 \leq H^{+} < 1000$~м~--- 2 балла, $H^{+} \geq 1000$~м~--- 3 балла.
\end{itemize}

Итоговый балл $S = S_D + S_{H^{+}}$ определяет категорию: $S \leq 2$~--- <<Лёгкий>>, $3 \leq S \leq 4$~--- <<Средний>>, $S \geq 5$~--- <<Сложный>>.

\subsection{Алгоритм сжатия изображений}

Для оптимизации хранения и передачи фотографий реализован алгоритм серверной обработки изображений, выполняемый асинхронно через брокер сообщений NATS. Алгоритм состоит из следующих этапов:

\begin{enumerate}
  \item Получение сообщения из очереди NATS с метаданными загруженной фотографии.
  \item Декодирование изображения из формата JPEG/PNG.
  \item Масштабирование: если ширина изображения превышает $W_{\max} = 1920$~пикселей, выполняется пропорциональное уменьшение с использованием фильтра Ланцоша (Lanczos3) для сохранения качества.
  \item Генерация миниатюры: масштабирование до ширины $W_{th} = 300$~пикселей для быстрой загрузки на карте.
  \item Кодирование в JPEG с качеством $Q = 85\%$~--- баланс между размером файла и визуальным качеством.
  \item Загрузка обработанных изображений в объектное хранилище MinIO (S3-совместимый интерфейс).
  \item Обновление записи в базе данных: сохранение адресов оригинала и миниатюры, установка статуса <<обработано>>.
\end{enumerate}
